\newpage
\chapter*{\Large NỘI DUNG}
\addcontentsline{toc}{chapter}{NỘI DUNG}
\chapter{KIẾN THỨC CHUẨN BỊ}

\section{Tập affine và tập lồi}
\label{section:affine_and_convex_sets}

\subsection{Đường thẳng và đoạn thẳng}
Giả sử $x_1 \neq x_2$ là hai điểm trong $\mathbb{R}^n$. Các điểm có dạng \(y = \theta x_1 + (1-\theta)x_2 \), trong đó $\theta \in \R$, tạo thành đường thẳng đi qua $x_1$ và $x_2$. Giá trị tham số $\theta = 0$ tương ứng với $y = x_2$, và giá trị tham số $\theta = 1$ tương ứng với $y = x_1$. Các giá trị của tham số $\theta$ nằm giữa 0 và 1 tương ứng với đoạn thẳng (đóng) giữa $x_1$ và $x_2$.

\noindent Biểu diễn $y$ dưới dạng
\begin{equation}
    y = x_2 + \theta(x_1 - x_2)
\end{equation}
cho ta một cách hiểu khác: $y$ là tổng của điểm cơ sở $x_2$ (tương ứng với $\theta = 0$) và vector chỉ phương $x_1 - x_2$ (chỉ từ $x_2$ đến $x_1$) được nhân với tham số $\theta$. Do đó, $\theta$ cho biết tỉ lệ khoảng cách từ $x_2$ đến $x_1$ mà điểm $y$ nằm trên đó. Khi $\theta$ tăng từ 0 đến 1, điểm $y$ di chuyển từ $x_2$ đến $x_1$; với $\theta > 1$, điểm $y$ nằm trên đường thẳng vượt qua $x_1$.

\subsection{Tập affine}
\begin{define}\textit{(Tập affine)}
    \label{define:affine_set}
    Một tập $\mathcal{C} \subseteq \R^n$ được gọi là tập affine nếu đường thẳng đi qua bất kỳ hai điểm phân biệt nào trong $\mathcal{C}$ đều nằm trong $\mathcal{C}$, nghĩa là với mọi $x_1, x_2 \in \mathcal{C}$ và $\theta \in \mathbb{R}$, ta có $\theta x_1 + (1-\theta)x_2 \in \mathcal{C}$. Nói cách khác, $\mathcal{C}$ chứa tổ hợp tuyến tính của bất kỳ hai điểm nào trong $\mathcal{C}$, miễn là các hệ số trong tổ hợp tuyến tính có tổng bằng một.
\end{define}

Ý tưởng này có thể được tổng quát hóa cho nhiều hơn hai điểm. Ta gọi một điểm có dạng
$\theta_1 x_1 + \cdots + \theta_k x_k$, trong đó $\theta_1 + \cdots + \theta_k = 1$, là tổ hợp affine của các điểm $x_1, \ldots, x_k$. Sử dụng quy nạp từ định nghĩa của tập affine (tức là tập chứa mọi tổ hợp affine của hai điểm bất kỳ trong nó), có thể chứng minh rằng tập affine chứa mọi tổ hợp affine của các điểm trong nó: nếu $\mathcal{C}$ là tập affine, $x_1, \ldots, x_k \in \mathcal{C}$, và $\theta_1 + \cdots + \theta_k = 1$, thì điểm $\theta_1 x_1 + \cdots + \theta_k x_k$ cũng thuộc $\mathcal{C}$.

\noindent Nếu $\mathcal{C}$ là tập affine và $x_0 \in \mathcal{C}$, thì tập
\begin{equation}
    \mathcal{V} = \mathcal{C} - x_0 = \{x - x_0 \mid x \in \mathcal{C}\}
\end{equation}
là một không gian con, tức là đóng với phép cộng và phép nhân với số vô hướng. Để thấy điều này, giả sử $v_1, v_2 \in \mathcal{V}$ và $\alpha, \beta \in \mathbb{R}$. Khi đó ta có $v_1 + x_0 \in \mathcal{C}$ và $v_2 + x_0 \in \mathcal{C}$, và do đó
\begin{equation}
    \alpha v_1 + \beta v_2 + x_0 = \alpha(v_1 + x_0) + \beta(v_2 + x_0) + (1-\alpha-\beta)x_0 \in \mathcal{C},
\end{equation}
vì $\mathcal{C}$ là tập affine và $\alpha + \beta + (1-\alpha-\beta) = 1$. Ta kết luận rằng $\alpha v_1 + \beta v_2 \in \mathcal{V}$ vì $\alpha v_1 + \beta v_2 + x_0 \in \mathcal{C}$.

\noindent Do đó, tập affine $\mathcal{C}$ có thể được biểu diễn dưới dạng
\begin{equation}
    \mathcal{C} = \mathcal{V} + x_0 = \{v + x_0 \mid v \in \mathcal{V}\},
\end{equation}
tức là một không gian con cộng với một điểm dịch chuyển. Không gian con $\mathcal{V}$ gắn với tập affine $\mathcal{C}$ không phụ thuộc vào việc chọn $x_0$, vì vậy $x_0$ có thể được chọn là bất kỳ điểm nào trong $\mathcal{C}$. Ta định nghĩa chiều của tập affine $\mathcal{C}$ là chiều của không gian con $\mathcal{V} = \mathcal{C} - x_0$, trong đó $x_0$ là một phần tử bất kỳ của $\mathcal{C}$.

Tập hợp tất cả các tổ hợp affine của các điểm trong một tập \( \mathcal{C} \subseteq \mathbb{R}^n \)
được gọi là \textit{bao affine} (affine hull) của \( \mathcal{C} \), và còn được ký hiệu là \( \operatorname{\textbf{aff}} \mathcal{C} \):
\[
\operatorname{\textbf{aff}} \mathcal{C} = \{ \theta_1 x_1 + \cdots + \theta_k x_k \mid x_1, \ldots, x_k \in \mathcal{C},\; \theta_1 + \cdots + \theta_k = 1 \}.
\]
Bao affine là tập affine nhỏ nhất chứa \( \mathcal{C} \), theo nghĩa sau:
nếu \( \mathcal{S} \) là bất kỳ tập affine nào thỏa \( \mathcal{C} \subseteq \mathcal{S} \),
thì \( \operatorname{\textbf{aff}} \mathcal{C} \subseteq \mathcal{S} \).

\subsection{Tập lồi}
\begin{define}\textit{(Tập lồi)}
    \label{define:convex_set}
    Một tập \( \mathcal{C} \) được gọi là \textit{tập lồi} nếu đoạn thẳng nối giữa hai điểm bất kỳ trong \( \mathcal{C} \) đều nằm trong \( \mathcal{C} \); tức là, với mọi \( x_1, x_2 \in \mathcal{C} \) và mọi \( \theta \) thỏa \( 0 \leq \theta \leq 1 \), ta có: \( \theta x_1 + (1 - \theta)x_2 \in \mathcal{C}\).
\end{define}

Mọi \textit{tập affine} cũng là tập lồi, bởi vì nó chứa toàn bộ đường thẳng đi qua hai điểm phân biệt bất kỳ trong nó, và do đó cũng chứa đoạn thẳng nối hai điểm đó.

\begin{define}\textit{(Tổ hợp lồi)}
    \label{define:convex_combination}
    Một điểm được gọi là \textit{tổ hợp lồi} (convex combination) của các điểm \( x_1, x_2, \ldots, x_k \) nếu nó có thể được viết dưới dạng
    \(
    x = \theta_1 x_1 + \theta_2 x_2 + \cdots + \theta_k x_k,
    \)
    với các hệ số \( \theta_1, \theta_2, \ldots, \theta_k \) thỏa mãn
    \[
    \theta_1 + \theta_2 + \cdots + \theta_k = 1, \quad \text{và} \quad \theta_i \geq 0, \; \forall i = 1, 2, \ldots, k.
    \]
\end{define}


Tương tự như trong trường hợp các tập affine, có thể chứng minh rằng một tập là lồi khi và chỉ khi nó chứa mọi tổ hợp lồi của các điểm của chính nó.
Một tổ hợp lồi có thể được hiểu như một hỗn hợp hay trung bình có trọng số
của các điểm, trong đó \( \theta_i \) là tỷ lệ (hay phần đóng góp) của \( x_i \) trong hỗn hợp.

\begin{define}\textit{(Bao lồi - Convex hull)}
    \label{define:convex_hull}
    \textit{Bao lồi} của một tập \( \mathcal{C} \), ký hiệu là \( \operatorname{\textbf{conv}} \mathcal{C} \), được định nghĩa là tập hợp tất cả các tổ hợp lồi của các điểm trong \( \mathcal{C} \):
    \begin{align*}
    \operatorname{\textbf{conv}} \mathcal{C}
    = \{\, \theta_1 x_1 + \cdots + \theta_k x_k
    \mid x_i \in \mathcal{C}, \; \theta_i \ge 0, \\ i = 1, \ldots, k, \;
    \theta_1 + \cdots + \theta_k = 1 \,\}.
    \end{align*}
\end{define}

Như tên gọi, \( \operatorname{\textbf{conv}} \mathcal{C} \) luôn là một tập lồi. Đây cũng là tập lồi nhỏ nhất chứa \( \mathcal{C} \), tức là, nếu \( \mathcal{B} \) là một tập lồi bất kỳ thỏa \( \mathcal{C} \subseteq \mathcal{B} \), thì \( \operatorname{\textbf{conv}} \mathcal{C} \subseteq \mathcal{B} \).

\subsection{Một số ví dụ về tập lồi}
\begin{define}\textit{(Siêu phẳng - Hyperplane)}
    \label{define:hyperplane}
    Một \textit{siêu phẳng} trong không gian \( n \)-chiều là tập hợp các điểm thỏa mãn phương trình
    \[
    a_1 x_1 + a_2 x_2 + \cdots + a_n x_n = \mb{a}^T \mb{x} = b,
    \]
    trong đó \( b, a_i \in \R^{n}, \; i = 1, 2, \ldots, n \).
\end{define}

\begin{define}\textit{(Nửa không gian - Halfspace)}
    \label{define:halfspace}
    Một \textit{nửa không gian} trong không gian \( n \)-chiều là tập hợp các điểm thỏa mãn bất phương trình
    \[
    a_1 x_1 + a_2 x_2 + \cdots + a_n x_n = \mb{a}^T \mb{x} \leq b,
    \]
    trong đó \( b, a_i \in \R^{n}, \; i = 1, 2, \ldots, n \).
\end{define}

\begin{define}\textit{(Hình cầu theo chuẩn - Norm ball)}
    \label{define:norm_ball}
    Cho một \textit{tâm} \( x_c \), một \textit{bán kính} \( r > 0 \), 
    và một \textit{chuẩn} (norm) dùng để xác định khoảng cách giữa các điểm. 
    \textit{Hình cầu theo chuẩn} tương ứng là tập hợp các điểm thỏa mãn:
    \[
    B(x_c, r) = \{\, x \mid \|x - x_c\|_2 \leq r \,\}
    = \{\, x_c + r u \mid \|u\|2 \leq 1 \,\}.
    \]
\end{define}

\begin{define}\textit{(Đa diện - Polyhedra)}
    \label{define:polyhedra}
    Một \textit{đa diện} được định nghĩa là tập nghiệm của một số hữu hạn các phương trình và bất phương trình tuyến tính:
    \[
    P = \{x \mid a_j^T x \le b_j, \; j = 1, \ldots, m, \; c_j^T x = d_j, \; j = 1, \ldots, p \}. \tag{2.5}
    \]
\end{define}

\noindent Có thể thấy một đa diện chính là giao của hữu hạn các nửa không gian và siêu phẳng. Các \textit{tập affine} (ví dụ: không gian con tuyến tính, siêu phẳng, đường thẳng), \textit{tia}, \textit{đoạn thẳng}, và \textit{nửa không gian} đều là các \textit{đa diện}.

\subsection{Các phép toán bảo toàn tính lồi}

\subsubsection{Phép giao}

Tính lồi được bảo toàn dưới phép giao, nếu \(\mathcal{S}_1\) và \(\mathcal{S}_2\) là các tập lồi, thì \(\mathcal{S}_1 \cap \mathcal{S}_2\) cũng là tập lồi. Tính chất này được mở rộng cho giao của vô hạn tập hợp:
nếu mỗi \(\mathcal{S}_\alpha\) là tập lồi với mọi \(\alpha \in A\), thì
\( \bigcap_{\alpha \in A} \mathcal{S}_\alpha \) là tập lồi (các \textit{không gian con tuyến tính}, \textit{tập affine}, cũng đều đóng dưới phép giao tùy ý).

\noindent Ví dụ, một \textit{đa diện} (polyhedron) là giao của các \textit{nửa không gian} (halfspace) và \textit{siêu phẳng} (hyperplane), vốn đều là các tập lồi, do đó đa diện cũng là một tập lồi.

\subsubsection{Hàm affine}

\begin{define}\textit{(Hàm affine)}
    \label{define:affine_func}
    Một hàm \( f : \mathbb{R}^n \to \mathbb{R}^m \) được gọi là 
    \textit{hàm affine} nếu nó là tổng của một \textit{hàm tuyến tính} và một \textit{hằng số}, tức là nó có dạng \( f(x) = A x + b, \) trong đó \(A \in \mathbb{R}^{m \times n}\) và \(b \in \mathbb{R}^m\).
\end{define}

\noindent Giả sử \(\mathcal{S} \subseteq \mathbb{R}^n\) là một \textit{tập lồi}, và \(f : \mathbb{R}^n \to \mathbb{R}^m\) là một hàm affine. Khi đó, \textit{ảnh của tập} \(\mathcal{S}\) \textit{qua ánh xạ} \(f\) được định nghĩa là
\[
f(\mathcal{S}) = \{ f(x) \mid x \in \mathcal{S} \}.
\]
là tập lồi. Tương tự, nếu \(f : \mathbb{R}^k \to \mathbb{R}^n\) là một hàm afine, \textit{nghịch ảnh} của tập \(\mathcal{S}\) qua ánh xạ \(f\) được định nghĩa là
\[
f^{-1}(\mathcal{S}) = \{ x \mid f(x) \in \mathcal{S} \}.
\]
cũng là \textit{tập lồi}.